% Source: https://artofproblemsolving.com/community/c3h2517157_why_does_this_feel_wrong
% Stable identity: aops-topic-2517157-q01
\textbf{(a)} Each face has probability $\dfrac16$, so
\[E(X)=\dfrac16(1+2+3+4+5+6)=\dfrac{21}{6}=\boxed{\dfrac72=3.5}.\]

\textbf{(b)} The total is $100(1+2+\cdots+6)=2100$, so the average is $\dfrac{2100}{600}=3.5$. In general, in a long run of rolls each face comes up about $\dfrac16$ of the time, so the average of the results is close to $E(X)$: the expected value is a long-run average.

\textbf{(c)} $\sum_{x=1}^{6}(x-3.5)=21-6\cdot3.5=0$: the deviations $-2.5,-1.5,-0.5,0.5,1.5,2.5$ cancel, so $3.5$ is the balance point of the six equally likely values. An average of whole numbers need not be a whole number (the average of $1$ and $2$ is $1.5$), so $E(X)$ need not be a possible value of $X$; it describes the centre of the distribution, not a single roll.
